% Table 3.4: Historical Macro-Climatic Yield Shocks in the Midwestern Panel (1951-2025)
\begin{table}[htbp]
\centering
\small
\caption{Historical macro-climatic shock years, farm-gate price responses, and associated drivers (1951--2025).}
\label{tab:yield_shocks}
\begin{tabularx}{\textwidth}{c c c c X}
\toprule
\textbf{Year} & \textbf{Yield Anom.} & \textbf{Rel. Anom.} & \textbf{Max Price Swing} & \textbf{Primary Meteorological / Ecological Driver} \\
 & (bu/ac) & (\%) & (\%) & \\
\midrule
\multicolumn{5}{l}{\textit{\textbf{Panel A: Major Adverse Macro-Climatic Yield Shocks (Negative Anomalies)}}} \\[0.5ex]
1988 & -9.64 & -25.6\% & +33.0\% & Historic severe summer drought and persistent heatwave across the Corn Belt. \\
2003 & -9.72 & -21.6\% & +59.5\% & Late-summer heat and moisture stress compounded by widespread soybean aphid outbreak. \\
1974 & -6.09 & -19.7\% & +58.6\% & Excessively wet spring delaying planting, followed by mid-summer drought and early frost. \\
2012 & -5.80 & -11.7\% & +17.4\% & Extensive summer flash drought and record July heat across the central United States. \\
1983 & -4.02 & -11.4\% & +36.0\% & Prolonged July-August heatwave and severe drought across the Corn and Soybean Belt. \\
1993 & -4.22 & -10.5\% & +17.5\% & Great Upper Mississippi River flood; prolonged waterlogging, root asphyxia, and disease. \\
\midrule
\multicolumn{5}{l}{\textit{\textbf{Panel B: Prominent Favorable Bumper Harvests (Positive Anomalies)}}} \\[0.5ex]
1952 & +2.41 & +11.9\% & --- & Exceptionally favorable vegetative and pod-filling moisture across the central belt. \\
2021 & +5.79 & +10.8\% & -14.4\% & Optimal reproductive rainfall and moderate temperatures throughout August. \\
2016 & +5.53 & +10.8\% & +13.2\% & Broadly distributed summer rains and minimal extreme heat stress across the Midwest. \\
1994 & +4.33 & +10.7\% & -19.3\% & Record-setting benign growing conditions with timely August rains and cool nights. \\
\bottomrule
\end{tabularx}
\vspace{1ex}
\raggedright
\footnotesize{\textit{Notes:} Yield Anom. is $y_t - \hat{y}_t$ and Rel. Anom. is $(y_t - \hat{y}_t)/\hat{y}_t \times 100$, where $\hat{y}_t = 0.484t - 925.46$ is estimated via linear OLS over the 135-county panel (1951--2025). Max Price Swing denotes the peak percentage departure in national farm-gate prices across a 12-month window centered on the crop campaign (April of year $t$ to March of year $t+1$), evaluated relative to the pre-planting baseline in April. For shortfall years, it captures the peak upward supply-rationing rally; for bumper harvests, it captures post-harvest cash price depression. \textit{Source:} Author's calculation on USDA NASS and farmdoc (University of Illinois) data.}
\end{table}
